\section{The round-sphere consequence}\label{sec:spheres}

\begin{proof}[Proof of \Cref{cor:round-spheres}]
For $(x,P)\in G_m(S^{m+n})$, define the cone by
\[
 \int\varphi\,\dd\mathcal C(V)
 =\int_{G_m(S^{m+n})}\int_0^\infty
   \varphi(rx,\R x\oplus P)r^m\,\dd r\,\dd V(x,P).
\]
The polar area formula \cite[Chapter~3, \S2, (2.4)--(2.6)]{Simon2018} gives
an integral $(m+1)$-varifold in $\R^{m+n+1}$ with the same integer
multiplicities away from the vertex. The vertex has zero mass, and
\[
 \|\mathcal C(V)\|(B_R(0))
 =\frac{R^{m+1}}{m+1}\|V\|(S^{m+n}).
\]

The radial and angular first-variation calculation below adapts the
hypersurface calculation in \cite[Lemma~6.1]{OpenAISolomonYau2026};
we include its derivation in the present dimensions.
For a compactly supported $C^1$ ambient field $X$ vanishing near the vertex,
write $X(rx)=f(r,x)x+Y_r(x)$ with $Y_r(x)\perp x$. Then
\[
 \Div_{\R x\oplus P}X(rx)
 =\partial_r f+\frac mr f+\frac1r\Div^S_PY_r.
\]
After multiplication by $r^m$, the radial terms integrate to the endpoint
values of $r^mf$, and the angular term vanishes by spherical stationarity
for every $r$. Thus the cone is stationary off the vertex. For a general
compactly supported field, a radial cutoff between radii $\eps$ and $2\eps$
has first-variation error $O_{X,V}(\eps^{m+1}+\eps^m)$ by the mass formula.
This tends to zero because $m\geq1$, so the cone is stationary everywhere.

The polar map identifies the punctured support of the cone with
$(0,\infty)\times\spt\|V\|$. A regular link patch gives a smooth embedded
local cone with the same constant multiplicity, and cone stationarity makes
it minimal. Conversely, suppose the cone is locally $q|M|$ near $rx$, with
$M$ smooth embedded and minimal. Its tangent planes contain the radial line
almost everywhere, hence everywhere by continuity. The radial vector field
is therefore tangent to $M$, and the radius function is a submersion on $M$.
The dilation flow identifies a smaller patch of $M$ with a radius interval
times a smooth embedded spherical section. Comparing the polar varifold
formulas on this product neighborhood identifies $V$ there with $q$ times
that section, including its tangent planes; spherical stationarity makes the
section minimal. Consequently
\[
 \operatorname{Sing}\mathcal C(V)\setminus\{0\}
 =\{rx:r>0,\ x\in\operatorname{Sing}_S V\}.
\]

Apply \Cref{thm:main} with dimensions $(m+1,n)$ to obtain
$\HH^{m+1}(\operatorname{Sing}\mathcal C(V))=0$. Let
$A=\operatorname{Sing}\mathcal C(V)\cap\{1<|z|<2\}$ and $\rho(z)=|z|$.
Regularity is local, so $A$ is Borel. The general Hausdorff slicing
inequality \cite[Chapter~2, \S1, Theorem~1.10]{Simon2018} applied to $\rho$
gives
\[
 \int_1^2\HH^m(A\cap\rho^{-1}(r))\,\dd r=0.
\]
For $1<r<2$, the slice is $r\operatorname{Sing}_S V$. Choosing a radius with
zero slice measure and using homogeneity gives $\HH^m(\operatorname{Sing}_S V)=0$ in the chordal
metric. The round and chordal metrics are locally bi-Lipschitz equivalent,
so they have the same Hausdorff null sets. Homotheties give the assertion for
every positive radius.
\end{proof}
